% Bagian: computability
% Bab: computability-theory
% Subbagian: prop-reduce

\documentclass[../../../include/open-logic-section]{subfiles}

\begin{document}

\olfileid[id]{cmp}{thy}{ppr}
\olsection{Sifat Redusibilitas}

Kita menulis $A \leq_m B$ jika $A$ tereduksi ke~$B$. Notasi ini menyarankan
bahwa jika $A \leq_m B$, maka $A$ ``tidak lebih sulit daripada''~$B$ dan $B$
``sama sulit atau lebih sulit daripada''~$A$, serta bahwa $\le_m$, seperti
$\le$ biasa pada bilangan, mengurutkan himpunan (atau masalah keputusan)
berdasarkan kompleksitasnya. Dua proposisi berikut mendukung intuisi ini.
Proposisi pertama menyatakan bahwa $\le_m$ bersifat transitif.

\begin{prop}
\ollabel{prop:trans-red}
Jika $A \leq_m B$ dan $B \leq_m C$, maka $A \leq_m C$.
\end{prop}

\begin{proof}
Mengomposisikan suatu reduksi dari $A$ ke~$B$ dengan suatu reduksi dari $B$ ke
$C$ menghasilkan suatu reduksi dari $A$ ke~$C$.
\end{proof}

\begin{prob}
Buktikan \olref{prop:trans-red} dengan menunjukkan bahwa jika $f$ dan~$g$
masing-masing merupakan reduksi banyak-ke-satu dari $A$ ke~$B$ dan dari $B$
ke~$C$, maka $\comp{f}{g}$ merupakan reduksi banyak-ke-satu dari $A$ ke~$C$.
\end{prob}

\begin{prop}
\ollabel{prop:reduce}
Misalkan $A$ dan $B$ sembarang himpunan, dan andaikan $A \leq_m B$.
\begin{enumerate}
\item Jika $B$ terenumerasi secara komputabel, maka demikian pula~$A$.
\item Jika $B$ dapat dihitung, maka demikian pula~$A$.
\end{enumerate}
\end{prop}

\begin{proof}
Misalkan $f$ suatu reduksi banyak-ke-satu dari $A$ ke~$B$. Untuk klaim
pertama, cukup periksa bahwa jika $B$ merupakan domain fungsi parsial~$g$,
maka $A$ merupakan domain~$\comp{f}{g}$:
\begin{align*}
x \in A & \text{ jika dan hanya jika } f(x) \in B \\
& \text{ jika dan hanya jika }  g(f(x)) \fdefined.
\end{align*}

Untuk klaim kedua, ingat bahwa jika $B$ dapat dihitung, maka $B$
dan~$\Complement{B}$ terenumerasi secara komputabel
(\olref[cmp]{thm:ce-comp}). Mudah diperiksa bahwa $f$ juga merupakan reduksi
banyak-ke-satu dari $\Complement{A}$ ke $\Complement{B}$. Jadi, berdasarkan
bagian pertama bukti ini, $A$ dan~$\Complement{A}$ keduanya
!!{computably enumerable}. Dengan demikian, $A$ juga dapat dihitung menurut
\olref[cmp]{thm:ce-comp}. (Sebagai alternatif, Anda dapat memeriksa bahwa
$\Char{A} = \comp{f}{\Char{B}}$; jadi jika $\Char{B}$ dapat dihitung, maka
demikian pula~$\Char{A}$.)
\end{proof}

\begin{prob}
Andaikan $f$ suatu reduksi banyak-ke-satu dari $A$ ke~$B$. Tunjukkan bahwa $f$
juga merupakan reduksi banyak-ke-satu dari $\Complement{A}$ ke
$\Complement{B}$.
\end{prob}

\begin{prob}
Tunjukkan bahwa jika $f\colon A \to B$ merupakan reduksi banyak-ke-satu, maka
$\Char{A} = \comp{f}{\Char{B}}$.
\end{prob}

\begin{digress}
Gagasan redusibilitas yang lebih umum, yang disebut \emph{redusibilitas
Turing}, berguna dalam konteks lain, khususnya untuk membuktikan hasil
ketakterputusan. Perhatikan bahwa menurut \olref[cmp]{cor:comp-k}, komplemen
$K_0$ tidak tereduksi ke~$K_0$ karena komplemen itu tidak terenumerasi secara
komputabel. Akan tetapi, secara intuitif, jika Anda mengetahui jawaban atas
pertanyaan tentang $K_0$, Anda juga mengetahui jawaban atas pertanyaan tentang
komplemennya. Suatu himpunan $A$ disebut tereduksi secara Turing ke~$B$ jika
jawaban atas pertanyaan tentang~$A$ dapat ditentukan dengan prosedur dapat
dihitung yang boleh mengajukan pertanyaan tentang~$B$. Gagasan ini lebih
longgar daripada redusibilitas banyak-ke-satu, yang hanya memperbolehkan Anda
(1)~mengajukan satu pertanyaan tentang $B$, dan (2)~menerjemahkan jawaban ``ya''
menjadi jawaban ``ya'' atas pertanyaan tentang $A$, dan demikian pula untuk
``tidak.'' Tetap benar bahwa jika $A$ tereduksi secara Turing ke~$B$ dan $B$
dapat dihitung, maka $A$ juga dapat dihitung (meskipun, seperti yang telah kita
lihat, pernyataan analognya tidak berlaku bagi enumerabilitas komputabel).

Anda sebaiknya memikirkan berbagai gagasan redusibilitas yang telah kita bahas
dan memahami perbedaannya. Namun, dalam bab ini kita hanya akan membahas
redusibilitas banyak-ke-satu. Kebetulan, kedua jenis redusibilitas yang dibahas
dalam paragraf terakhir mempunyai analog dalam kompleksitas komputasional,
dengan syarat tambahan bahwa mesin Turing berjalan dalam waktu polinomial:
versi kompleksitas redusibilitas banyak-ke-satu dikenal sebagai
\emph{redusibilitas Karp}, sedangkan versi kompleksitas redusibilitas Turing
dikenal sebagai \emph{redusibilitas Cook}.
\end{digress}

\end{document}

